\chapter*{Abstract}
\addcontentsline{toc}{chapter}{Abstract}

Federated Learning has established itself as a foundational paradigm for privacy-preserving collaborative model training across decentralized data sources. However, scaling modern foundation models to billions of parameters exposes a critical dual bottleneck in federated optimization: an excessive \emph{spatial} communication overhead from transmitting high-dimensional parameter updates, and an inefficient \emph{temporal} communication pattern imposed by rigid, fixed-interval synchronization schedules. Parameter-efficient fine-tuning methods such as Low-Rank Adaptation address the spatial dimension by decomposing updates into compact low-rank factors. Yet, naively applying Low-Rank Adaptation in federated environments introduces a fundamental structural inconsistency: independent averaging of constituent low-rank matrices fails to reconstruct the true global update due to the non-commutativity of matrix multiplication. Furthermore, existing federated Low-Rank Adaptation variants remain temporally inefficient by relying on static synchronization frequencies agnostic to local training dynamics.

To address both dimensions simultaneously, this thesis investigates combining Federated Dynamic Averaging (FDA) with federated low-rank adaptation. Specifically, we develop and evaluate the FDA-RoLoRA Dual Track protocol, an algorithmic formulation that pairs dynamic, variance-driven averaging with Robust Low-Rank Adaptation (RoLoRA). This method eliminates the bilinear aggregation inconsistency through alternating updates between down-projection and up-projection matrices, while dynamically regulating synchronization frequency via decoupled, smoothed per-matrix variance thresholds ($\theta_A, \theta_B$). While full-scale evaluations on GLUE benchmarks using RoBERTa-Large and Llama-2-7B are in progress at the time of writing, preliminary results are reported on a controlled synthetic benchmark under non-IID data distributions against standard FedAvg-LoRA, fixed-schedule RoLoRA, and FFA-LoRA baselines to quantify communication efficiency, convergence stability, and local computation trade-offs.



